\scanpage{38}
6) $\Gamma\simeq\mathcal A_5$.

\[
\begin{aligned}
Z(\Gamma)&=\{1\}\\ 
N(\Gamma)&=\Gamma\simeq\mathcal A_5
\end{aligned}
\]

\uuline{Une fois} choisi un cas de référence \inserted{$(G_0,\Gamma_0)$} sur $\mathbb Q$, si $S$ est
un schéma de car. 0, se donner un couple
$(G\supset\Gamma)$, avec $G$ forme de $\mathrm{GP}(1)$, $\Gamma$ forme de $\mathcal A_5$,
revient au $=$ que de se donner une forme
\uline{intérieure} de $\Gamma_0$, ou un torseur sous $\Gamma_0$.
Si on \deleted{avait} \inserted{choisit} $\Gamma_{0k}=\mathcal A_5$, il revient au $=$ de se donner
\deleted{\unsure{un}} un rev. de degré 5 de $S$ + trivialisation
du rev. quadratique de ses orientations.

Mais on a ceci : $\exists$ \inserted{sur $k$} forme non tordue de $\mathcal A_5$ ds un
$G$ (forme de $\mathrm{GP}(1)$) sss l'équation $\frac12$ cyclotomique
$\alpha^2+\alpha-1=0$ a ses racines dans $k$. \deleted{et alors il existe}
\deleted{une fo} \inserted{Pour que de} plus $G$ soit non tordu, il f et s. que
$k$ contienne les racines cinquièmes de $1$.

% slanted
\margin{\uuline{vérifier}}

Sur un $k$ \uline{muni} d'une solution $\alpha$, il y a un
plongmt de $\mathcal A_5$ dans un $G$ [défini à iso unique
près \inserted{\deleted{\unsure{de la Tr}} en exigeant que la valeur de $\unsure{c|A|\mathrm{Tr}}$ dfie par une permutation circulaire standart soit cet $\alpha$ !}]
-- OPS $k=k_0=\mathbb Q[T]/(T^2+T-1)\simeq\mathbb Q[\alpha_0]$, où
$\alpha_0=\frac12(-1\mathbin{\oplus}\sqrt5)$] Sur les $k_0$-schémas, la donnée
d'un $G\supset\Gamma$, $\Gamma$ \inserted{($G$)} forme de $\mathcal A_5$ ($\mathrm{GP}(1)$) équivaut
alors à la donnée d'un torseur sous $\mathcal A_5$,
i.e. ---

% archive sheet number: 37
